\documentclass[a4paper,12pt]{article}
\usepackage{amsmath, amssymb}

\begin{document}
	
	\section*{Laplace Transform and Bode Plots}
	
	The Laplace transform of a causal function $f(t)$ is defined as
	
	$$
	F(s) = \int_{0}^{\infty} f(t)\, e^{-st}\, dt, 
	\qquad s = \sigma + j\omega.
	$$
	
	Here the complex variable $s$ has two components:
	\begin{itemize}
		\item $\sigma$ (real part) controls exponential growth or decay.
		\item $j\omega$ (imaginary part) represents oscillation at frequency $\omega$.
	\end{itemize}
	
	\subsection*{Restriction to the Imaginary Axis}
	
	If we are interested only in the steady--state sinusoidal response of a system, 
	we set $\sigma = 0$. Thus
	
	$$
	s = j\omega,
	$$
	
	and the Laplace transform reduces to
	
	$$
	F(j\omega) = \int_{0}^{\infty} f(t)\, e^{-j\omega t}\, dt,
	$$
	
	which is exactly the Fourier transform of $f(t)$.  
	This corresponds to evaluating the Laplace transform along the \emph{imaginary axis} of the $s$--plane.
	
	\subsection*{Transfer Functions and Frequency Response}
	
	For a linear time--invariant (LTI) system with input $x(t)$ and output $y(t)$, 
	the transfer function is defined as
	
	$$
	H(s) = \frac{Y(s)}{X(s)}.
	$$
	
	To study the frequency response, we evaluate
	
	$$
	H(j\omega).
	$$
	
	The two quantities of interest for Bode analysis are:
	\begin{itemize}
		\item Magnitude: 
		$$
		|H(j\omega)|,
		$$
		which gives amplitude scaling.
		\item Phase: 
		$$
		\arg \big( H(j\omega) \big),
		$$
		which gives phase shift.
	\end{itemize}
	
	\subsection*{Bode Plots}
	
	The Bode plots are defined as
	$$
	20 \log_{10} |H(j\omega)| \quad \text{(dB magnitude plot)},
	$$
	and
	$$
	\arg \big( H(j\omega) \big) \quad \text{(phase plot)}.
	$$
	
	Thus, Bode plots are simply the magnitude and phase of the transfer function 
	evaluated on the imaginary axis of the Laplace domain.
	
	\subsection*{Interpretation}
	
	\begin{itemize}
		\item The \textbf{Laplace domain} ($s$--plane) describes both transient and steady--state behavior.
		\item Restricting to the line $s=j\omega$ isolates \textbf{pure frequency response}.
		\item This restriction is mathematically a \emph{slice} of the $s$--plane, corresponding to the Fourier transform.
	\end{itemize}
	
	\section*{Example: First-Order RC Low-Pass Filter}
	
	\subsection*{1. Circuit Equation (Time Domain)}
	
	For a resistor $R$ and capacitor $C$ in series with input $x(t)$ applied across both
	and output $y(t)$ taken across the capacitor, we have
	
	$$
	v_{in}(t) = v_{R}(t) + v_{C}(t).
	$$
	
	Using Ohm’s law and the capacitor law,
	$$
	v_{R}(t) = R i(t), 
	\qquad 
	i(t) = C \frac{dv_{C}(t)}{dt}.
	$$
	
	Thus
	$$
	v_{in}(t) = RC \frac{dv_{C}(t)}{dt} + v_{C}(t).
	$$
	
	\subsection*{2. Input--Output Relation}
	
	Let $x(t) = v_{in}(t)$ and $y(t) = v_{C}(t)$. Then
	$$
	x(t) = RC \frac{dy(t)}{dt} + y(t).
	$$
	
	\subsection*{3. Laplace Transform}
	
	Applying the Laplace transform (assuming zero initial conditions):
	$$
	X(s) = RC \, s Y(s) + Y(s).
	$$
	
	\subsection*{4. Transfer Function}
	
	The transfer function is
	$$
	H(s) = \frac{Y(s)}{X(s)} = \frac{1}{1 + RCs}.
	$$
	
	\subsection*{5. Frequency Response}
	
	Evaluating on the imaginary axis $s=j\omega$:
	$$
	H(j\omega) = \frac{1}{1 + j \omega RC}.
	$$
	
	Magnitude:
	$$
	|H(j\omega)| = \frac{1}{\sqrt{1 + (\omega RC)^2}}.
	$$
	
	Phase:
	$$
	\arg \big( H(j\omega) \big) = -\tan^{-1}(\omega RC).
	$$
	
	\subsection*{6. Bode Plot Interpretation}
	
	\begin{itemize}
		\item Low frequencies ($\omega \ll 1/RC$): 
		$|H(j\omega)| \approx 1$, phase $\approx 0^\circ$. 
		\hfill (Passes low frequencies)
		\item High frequencies ($\omega \gg 1/RC$): 
		$|H(j\omega)| \approx 1/(\omega RC)$, phase $\approx -90^\circ$. 
		\hfill (Attenuates high frequencies)
	\end{itemize}
	
	The cutoff frequency (at $-3$ dB) is
	$$
	\omega_c = \frac{1}{RC}.
	$$
	
	\section*{Example: Mechanical Low-Pass (Spring--Damper in Parallel)}
	
	\subsection*{Analogy}
	Electrical $\leftrightarrow$ Mechanical correspondence:
	$$
	\begin{aligned}
		\text{Voltage} &\;\leftrightarrow\; \text{Force}, \\
		\text{Current} &\;\leftrightarrow\; \text{Velocity}, \\
		\text{Resistor }(R) &\;\leftrightarrow\; \text{Damper }(b), \\
		\text{Capacitor }(C) &\;\leftrightarrow\; \text{Spring }(k).
	\end{aligned}
	$$
	
	
	\subsection*{1. Governing Equation (Time Domain)}
	A damper (coefficient $b$) and spring (stiffness $k$) in parallel are driven by an external force input $x(t)$.
	Let the output be the displacement $y(t)$ (same for both parallel elements). Force balance gives
	$$
	x(t) = b\,\frac{dy(t)}{dt} + k\,y(t).
	$$
	
	\subsection*{2. Laplace Transform}
	Assuming zero initial conditions,
	$$
	X(s) = b\,s\,Y(s) + k\,Y(s).
	$$
	
	\subsection*{3. Transfer Function}
	$$
	H(s) \;=\; \frac{Y(s)}{X(s)} \;=\; \frac{1}{b s + k}.
	$$
	
	\subsection*{4. Frequency Response}
	Evaluating on the imaginary axis $s=j\omega$:
	$$
	H(j\omega) = \frac{1}{k + j\omega b}.
	$$
	Magnitude:
	$$
	\big|H(j\omega)\big| = \frac{1}{\sqrt{k^2 + (\omega b)^2}}.
	$$
	Phase:
	$$
	\arg\!\big(H(j\omega)\big) = -\tan^{-1}\!\left(\frac{\omega b}{k}\right).
	$$
	
	\subsection*{5. Bode Interpretation and Cutoff}
	Low frequency ($\omega \ll k/b$): spring dominates, $\;|H(j\omega)| \approx 1/k,\; \arg \approx 0^\circ.$
	
	High frequency ($\omega \gg k/b$): damper dominates, $\;|H(j\omega)| \approx 1/(\omega b),\; \arg \approx -90^\circ.$
	
	Cutoff (–3 dB) frequency:
	$$
	\omega_c = \frac{k}{b}.
	$$
	
	\subsection*{6. Parameter Mapping to RC Filter}
	Matching the RC form $H(s)=\dfrac{1}{1 + RC\,s}$ via factorization,
	$$
	H(s) = \frac{1}{k}\,\frac{1}{1 + \left(\frac{b}{k}\right) s}
	\quad\Longleftrightarrow\quad
	\text{gain } \frac{1}{k},\;\; \text{time constant } \tau = \frac{b}{k}
	\;\;(\text{cf. } \tau_{\text{RC}}=RC).
	$$
	
	
\end{document}
